\documentclass{article}

% Math support
\usepackage{amsmath}
\usepackage{amssymb}

% For \citet, \citep, \citealp and bibliography
\usepackage{natbib}

% Nice hyperlinks (optional but harmless)
\usepackage{hyperref}

\begin{document}

\subsection{Kappa and weighting estimators of the LATE}
The previous section established Abadie's kappa representation of the LATE. For
the empirical analysis, the relevant implication is that this representation
identifies the population parameter, but does not by itself select a unique
finite-sample estimator. In population, the complier share
$P(D_i(1)>D_i(0))$ can be written as $E[\kappa_i]$, $E[\kappa_{1i}]$, or
$E[\kappa_{0i}]$. These objects are equal under the IV assumptions, but their
sample analogues need not be equal. As a result, the same identification result
leads to several natural weighting estimators once the population moments are
replaced by sample averages.
Following \citet{sloczynski2025abadie}, I focus on five such estimators. The
distinction that matters for the later analysis is whether the weights are
normalized in the sample. By normalization I mean that the weights entering a
weighted mean are rescaled to sum to one. For a treatment-effect contrast, this
requires the two weighted components of the contrast to be normalized separately.
The Uysal estimator $\hat\tau_u$ and the Abadie--Cattaneo estimator
$\hat\tau_{a,10}$ have this property. The other three estimators,
$\hat\tau_a$, $\hat\tau_t(=\hat\tau_{a,1})$, and $\hat\tau_{a,0}$, are
unnormalized sample analogues. This subsection defines the estimators and sets up
the notation used in the empirical tables. The implications of the normalization
choice for translation invariance are discussed in the next subsection.

\subsubsection*{Two normalized estimators}
For notational simplicity, the formulas below are written in terms of $p(X_i)$. In the empirical implementation, $p(X_i)$ is replaced by a fitted propensity score
$\hat p(X_i)$.
The first estimator, $\hat\tau_u$, is due to \citet{uysal2011three}. It is
the ratio of two separately Hájek-normalized inverse-probability-weighted
contrasts,
\begin{equation}
    \hat\tau_u =
    \frac{
        \left(\sum_{i=1}^N \frac{Z_i}{p(X_i)}\right)^{-1}\sum_{i=1}^N \frac{Y_iZ_i}{p(X_i)}
        -
        \left(\sum_{i=1}^N \frac{1-Z_i}{1-p(X_i)}\right)^{-1}\sum_{i=1}^N \frac{Y_i(1-Z_i)}{1-p(X_i)}
    }{
        \left(\sum_{i=1}^N \frac{Z_i}{p(X_i)}\right)^{-1}\sum_{i=1}^N \frac{D_iZ_i}{p(X_i)}
        -
        \left(\sum_{i=1}^N \frac{1-Z_i}{1-p(X_i)}\right)^{-1}\sum_{i=1}^N \frac{D_i(1-Z_i)}{1-p(X_i)}
    }.
    \label{eq:tau_u}
\end{equation}
Both the numerator and the denominator are differences of Hájek-normalized
IPW means. Thus, within each instrument group, the weights are rescaled before
the difference is formed.\\
The second normalized estimator, $\hat\tau_{a,10}$, is due to
\citet{abadie2018econometric} (see also \citealp{sant2022covariate}). Rather than normalizing the instrument-group IPW means as in \(\hat\tau_u\),
it normalizes the \(\kappa_1\)- and \(\kappa_0\)-weighted outcome means directly,
\begin{equation}
    \hat\tau_{a,10} =
    \left(\sum_{i=1}^N \kappa_{1i}\right)^{-1}\sum_{i=1}^N \kappa_{1i}Y_i
    -
    \left(\sum_{i=1}^N \kappa_{0i}\right)^{-1}\sum_{i=1}^N \kappa_{0i}Y_i.
    \label{eq:tau_a10}
\end{equation}

\subsubsection*{Three unnormalized estimators}
The remaining estimators start from the same population identity,
\[
\tau_{LATE}
= P(D_i(1)>D_i(0))^{-1}
E\!\left[
Y_i\frac{Z_i-p(X_i)}{p(X_i)(1-p(X_i))}
\right].
\]
They use the same sample numerator. What differs across them is the sample
analogue used for the complier share in the denominator.
\begin{align}
    \hat\tau_a &= \left(N^{-1}\sum_{i=1}^N \kappa_i\right)^{-1} N^{-1}\sum_{i=1}^N Y_i\,\frac{Z_i-p(X_i)}{p(X_i)(1-p(X_i))}, \label{eq:tau_a}\\
    \hat\tau_t \;(=\hat\tau_{a,1}) &= \left(N^{-1}\sum_{i=1}^N \kappa_{1i}\right)^{-1} N^{-1}\sum_{i=1}^N Y_i\,\frac{Z_i-p(X_i)}{p(X_i)(1-p(X_i))}, \label{eq:tau_t}\\
    \hat\tau_{a,0} &= \left(N^{-1}\sum_{i=1}^N \kappa_{0i}\right)^{-1} N^{-1}\sum_{i=1}^N Y_i\,\frac{Z_i-p(X_i)}{p(X_i)(1-p(X_i))}. \label{eq:tau_a0}
\end{align}
These three estimators are consistent under the same population argument, but
their finite-sample weights are not normalized. In particular, the two components
of the implied outcome contrast are not both forced to sum to one. This
distinction has direct finite-sample consequences. \citet{sloczynski2025abadie}
show that
$\hat\tau_u$ and $\hat\tau_{a,10}$ are translation invariant, whereas
$\hat\tau_a$, $\hat\tau_t(=\hat\tau_{a,1})$, and $\hat\tau_{a,0}$ are not.
Among the unnormalized estimators, $\hat\tau_t$ is the one most commonly used in
the applied literature, following \citet{tan2006regression} and
\citet{frolich2007nonparametric} and later applications such as
\citet{macurdy2011flexible},
\citet{donald2014inverse,donald2014testing}, and
\citet{abdulkadirouglu2017research}.

\subsubsection*{Estimating the propensity score}

In applications, the instrument propensity score $p(X_i)$ is usually unknown.
Following \citet{sloczynski2025abadie}, I use the parametric logit model
\[
F(X_i,\alpha)=\frac{\exp(X_i'\alpha)}{1+\exp(X_i'\alpha)}.
\]
I use two estimators of $\alpha$. The first is standard maximum likelihood,
denoted $\hat\alpha^{ml}$. The second is the covariate balancing propensity
score estimator of \citet{imai2014covariate}, denoted $\hat\alpha^{cb}$, which
solves
\begin{equation}
    N^{-1}\sum_{i=1}^N \frac{Z_i}{F(X_i,\hat\alpha^{cb})}X_i =
    N^{-1}\sum_{i=1}^N \frac{1-Z_i}{1-F(X_i,\hat\alpha^{cb})}X_i.
    \label{eq:cbps_moment}
\end{equation}
This moment condition chooses the propensity-score parameters so that the
reweighted covariates balance across instrument values. It therefore targets a
different finite-sample property than the likelihood score.
The superscript on an estimator indicates which propensity score is used. For example,
$\hat\tau_u^{ml}$ and $\hat\tau_u^{cb}$ are obtained by plugging in
$\hat p^{ml}(X_i)=F(X_i,\hat\alpha^{ml})$ and
$\hat p^{cb}(X_i)=F(X_i,\hat\alpha^{cb})$, respectively. The same convention applies
to the remaining kappa estimators.

This convention matters in the empirical analysis. Under maximum likelihood, the
finite-sample denominator ambiguity remains, and the ML versions of the
estimators generally differ. Under covariate balancing with a constant included
in $X_i$, \citet{sloczynski2025abadie} show the exact finite-sample identity
\[
\hat\tau_u^{cb}
=
\hat\tau_t^{cb}
=
\hat\tau_{a,1}^{cb}
=
\hat\tau_{a,0}^{cb}
=
\hat\tau_{a,10}^{cb}.
\]
Thus, reporting all these CBPS-based estimates separately would not add
information, because they coincide by construction in the sample. I therefore report
$\hat\tau_u^{cb}$ as the CBPS-based representative. The remaining estimators are
reported with maximum-likelihood propensity scores. The estimator \(\hat\tau_a^{cb}\) is not covered by this equivalence. Since the
empirical comparison follows \citet{sloczynski2025abadie} and focuses on
\(\hat\tau_u^{cb}\) as the CBPS-based representative, I do not report
\(\hat\tau_a^{cb}\) as a separate benchmark.







\subsection{Translation invariance}
\label{subsec:invariance_normalization}

The preceding subsection introduced several weighting estimators that are all
motivated by the same LATE identification argument. Their empirical behavior can
nevertheless differ sharply because they do not impose the same finite-sample
normalization. The most relevant consequence for this thesis is translation
invariance. An estimator is translation invariant if adding a constant to every
outcome leaves the estimated treatment effect unchanged. With \(Y\) denoting the
outcome vector and \(W\) all non-outcome variables used by the estimator, this
property can be written as
\begin{equation}
    \hat\tau(Y,W)=\hat\tau(Y+c\mathbf{1}_N,W)
    \quad\text{for all } c\in\mathbb{R}.
    \label{eq:ti_definition}
\end{equation}
\citet{sloczynski2025abadie} use this criterion to separate the normalized and
unnormalized kappa estimators. Their Proposition 3.2 shows that
\(\hat\tau_u\) and \(\hat\tau_{a,10}\) are translation invariant, whereas
\(\hat\tau_a\), \(\hat\tau_t(=\hat\tau_{a,1})\), and \(\hat\tau_{a,0}\) are
not. This result gives normalization a concrete interpretation: it is not only a
descriptive property of the weights, but determines whether the estimate depends
on an arbitrary additive coding choice for the outcome.

The same issue appears as scale sensitivity when the outcome is logarithmic. If
an income or wage variable \(V_i\) is measured in cents rather than dollars before
taking logs, the transformed outcome changes by a constant:
\begin{equation}
    \log(100V_i)=\log(V_i)+\log(100).
    \label{eq:log_units}
\end{equation}
An estimator that is not translation invariant in log outcomes can therefore
change when the underlying monetary variable is recorded in dollars, cents, or
larger units. This is the logic behind the cents-versus-dollars comparison in
\citet{sloczynski2025abadie}. The same argument also applies to binary outcomes
when one considers additive recodings such as \(Y_i+1\). By contrast, changing
from \(Y_i\) to \(1-Y_i\) is not a pure translation, because it combines a shift
with a sign reversal; it is therefore a different diagnostic question.

More generally, \citet{sloczynski2025abadie} formulate scale equivariance as a
requirement on how an estimator reacts to transformations of the outcome scale.
For the present analysis, the relevant special case is scale invariance of log
outcomes with respect to the units of the underlying variable. Multiplying the
untransformed outcome by a positive constant only adds a constant after the log
transformation. Hence, in this setting, the practical scale problem is captured
by the translation-invariance condition in \eqref{eq:ti_definition}.

The outcome-weight representation makes the translation-invariance condition
transparent. Suppose that, for a given fitted estimator, the point estimate can be
written as a weighted sum of observed outcomes,
\begin{equation}
    \hat\tau(Y)=\sum_{i=1}^N \omega_iY_i=\omega'Y,
    \label{eq:ow_rep_ch2}
\end{equation}
where \(\omega_i\) are the final outcome weights. These are not the same object
as Abadie's \(\kappa\)-weights. The \(\kappa\)-weights identify complier moments;
the outcome weights describe how a particular finite-sample estimator combines
the observed outcomes after all nuisance estimation and implementation choices
have been made. This distinction is central in \citet{knaus2024treatment}, who
shows how outcome weights can be derived and used to inspect DML and generalized
random forest estimators.

For a fixed vector of outcome weights, shifting all outcomes by \(c\) gives
\begin{align}
    \hat\tau(Y+c\mathbf{1}_N)
    &= \sum_{i=1}^N \omega_i(Y_i+c) \nonumber\\
    &= \hat\tau(Y)+c\sum_{i=1}^N\omega_i .
    \label{eq:ti_shift}
\end{align}
Thus, within a fixed weighted representation,
\begin{equation}
    \hat\tau(Y+c\mathbf{1}_N)=\hat\tau(Y)
    \quad\Longleftrightarrow\quad
    \sum_{i=1}^N\omega_i=0.
    \label{eq:ti_sumzero_ch2}
\end{equation}
The total weight sum is therefore the key finite-sample diagnostic. If
\(\sum_i\omega_i\neq0\), an additive shift of size \(c\) changes the estimate by
\(c\sum_i\omega_i\). For the log unit change in \eqref{eq:log_units}, the
predicted change is \(\log(100)\sum_i\omega_i\).

This condition is weaker than full normalization of treated and untreated outcome
weights. In a binary-treatment setting, full normalization requires
\begin{equation}
    \sum_{i=1}^N\omega_iD_i=1
    \quad\text{and}\quad
    \sum_{i=1}^N\omega_i(1-D_i)=-1.
    \label{eq:full_normalization}
\end{equation}
Full normalization implies \(\sum_i\omega_i=0\), and therefore implies
translation invariance. The reverse is not generally true: an estimator can have
total weight zero while the treated and untreated components do not separately
sum to \(1\) and \(-1\). Knaus refers to this weaker case as scale-normalization.
This distinction matters for the later diagnostics. Translation invariance asks
whether the estimate is stable to additive outcome recodings; full normalization
additionally asks whether the outcome-weight representation has the stronger
interpretation of a normalized treated mean minus a normalized untreated mean.

The empirical analysis uses two complementary checks. The first is an algebraic
frozen-weight check. After fitting the estimator once on an original outcome
\(Y\), I extract the outcome weights and keep them fixed. For the shifted outcome
\(Y+c\mathbf{1}_N\), the predicted shift is
\begin{equation}
    \Delta_B(c)
    =
    \sum_{i=1}^N\omega_i(Y_i+c)-\sum_{i=1}^N\omega_iY_i
    =
    c\sum_{i=1}^N\omega_i .
    \label{eq:method_b}
\end{equation}
This is the Method-B diagnostic in the empirical section. It is exact for the
extracted weights and directly measures the sensitivity implied by the fitted
outcome-weight representation.

The second check reruns the full estimation pipeline on the shifted outcome:
\begin{equation}
    \Delta_A(c)
    =
    \hat\tau(Y+c\mathbf{1}_N,W)-\hat\tau(Y,W).
    \label{eq:method_a}
\end{equation}
This is Method A. It is the main diagnostic in the later analysis because it asks
whether the estimator as actually implemented is invariant, including nuisance
training, cross-fitting, and tuning. For simple kappa estimators, the weights do
not depend on the outcome, so Method A and Method B coincide up to numerical
precision. For DML and forest-based estimators, the distinction is more
substantive. \citet{knaus2024treatment} shows that the availability and
properties of outcome weights depend on implementation choices. Standard
implementations of PLR-IV, instrumental forests, and Wald-AIPW are typically
scale-normalized rather than fully normalized; boosted-tree learners may also
fail the affine-smoother conditions used to obtain clean weighted
representations. Consequently, Method B is useful as a weight-based diagnostic,
but Method A is the more direct check for the estimator that is actually
reported.

Because Method A reruns an estimator, algorithmic randomness has to be controlled.
The same sample, random seed, cross-fitting folds, learner class, and tuning rule
are kept fixed across the original and shifted outcomes. This is especially
important for random forests, Ranger, and XGBoost, where sample splitting, tree
construction, and nested tuning can otherwise introduce unrelated variation. The
purpose of the rerun check is not to compare two independently randomized fits,
but to isolate the effect of the additive outcome shift.
\end{document}
